% TOP_PROBLEM: {"id": 30006683, "title": "Unconditional Separation of Quantum and Classical Bounded-Error Polynomial Time (BQP != BPP)", "category_id": 15, "category": {"id": 15, "name": "computer_science", "display_name": "Computer Science", "description": "Computational complexity, algorithms, and theoretical CS.", "slug": "computer-science", "order_index": 15, "created_at": "2026-07-31T15:26:25.671Z"}, "status": "open"}
% FIRST_POSTED: null
% !TeX program = lualatex
% Compile twice: lualatex open_problems_500.tex
% All references and prose are embedded; no BibTeX database or external inputs.
\documentclass[11pt,a4paper]{article}
\usepackage[margin=24mm,headheight=14pt]{geometry}
\usepackage{fontspec}
\setmainfont{Latin Modern Roman}
\setsansfont{Latin Modern Sans}
\setmonofont{Latin Modern Mono}
\usepackage{amsmath,amssymb,mathtools}
\usepackage{unicode-math}
\setmathfont{Latin Modern Math}
\usepackage{microtype}
\usepackage{enumitem,xurl,needspace}
\usepackage[dvipsnames]{xcolor}
\usepackage{hyperref}
\hypersetup{unicode=true,colorlinks=true,linkcolor=MidnightBlue,urlcolor=MidnightBlue,
 pdftitle={Mathematical Research Notebook},pdfauthor={Source-annotated compilation},bookmarksnumbered=true}
\usepackage{bookmark}
\usepackage{tocloft}
\setlength{\cftsecnumwidth}{3em}
\cftsetpnumwidth{2.5em}
\cftsetrmarg{4em}
\usepackage{titlesec}
\titleformat{\section}{\Large\bfseries\raggedright}{\thesection}{0.7em}{}
\usepackage{fancyhdr}
\pagestyle{fancy}\fancyhf{}
\fancyhead[L]{\small\sffamily Open Problems Math | Research notebook}
\fancyhead[R]{\small 27 September 2026}
\fancyfoot[C]{\thepage}
\setlength{\parindent}{0pt}
\setlength{\parskip}{4pt plus 1pt}
\setlength{\emergencystretch}{3em}
\setcounter{tocdepth}{1}
\setcounter{secnumdepth}{1}
\setlist[itemize]{leftmargin=3em,itemsep=3pt,topsep=4pt,parsep=0pt}
\allowdisplaybreaks
\newcommand{\N}{\mathbb{N}}
\newcommand{\Z}{\mathbb{Z}}
\newcommand{\Q}{\mathbb{Q}}
\newcommand{\R}{\mathbb{R}}
\newcommand{\C}{\mathbb{C}}
\newcommand{\F}{\mathbb{F}}
\newcommand{\A}{\mathbb{A}}
\newcommand{\PP}{\mathbb P}
\newcommand{\E}{\mathbb{E}}
\newcommand{\eps}{\varepsilon}
\newcommand{\dd}{\,\mathrm d}
\newcommand{\sref}[2]{\hyperlink{source:#1:#2}{[S#2]}}
\newcommand{\eref}[2]{\hyperlink{extra:#1:#2}{[E#2]}}
% Turn the next switch off for a shorter reading copy. The quotations preserve
% every exactTarget field, including qualifications omitted from a short summary.
\newif\ifshowcatalogue
\showcataloguetrue
\newcommand{\cataloguescope}[1]{\ifshowcatalogue
 \paragraph{Exact scope in the supplied catalogue.}
 \begin{quote}\small #1\end{quote}\fi}
\usepackage{etoolbox}
\AtBeginEnvironment{itemize}{\small}
% Standalone entry; original section content is delimited for verification.
\begin{document}
\setcounter{section}{0}
%% BEGIN ORIGINAL SECTION CONTAINER
% ================================================================
% problemId: 30006683
% Canonical problem ID: 30006683
% exactTarget SHA-256: bd8f49888411b74e177dcf1373a35d9d54729ac60bbdf7eae36f1af03fa4a889
\clearpage
\section*{\texorpdfstring{Unconditional Separation of Quantum and Classical Bounded-Error Polynomial Time (BQP != BPP)}{Unconditional Separation of Quantum and Classical Bounded-Error Polynomial Time (BQP != BPP)}}\label{30006683}
\subsection{Definitions and mathematical statement}
Both classes decide total binary languages with error at most $1/3$ on every input and polynomial running time. The class $\mathsf{BPP}$ uses classical random bits; $\mathsf{BQP}$ uses uniform quantum computation over a fixed effective universal gate set. The target is
\[
 \exists L\subseteq\{0,1\}^*:\quad L\in\mathsf{BQP}\ \text{and}\ L\notin\mathsf{BPP}.
\]
The assertion is unconditional and unrelativized: a separation only in the presence of an oracle or under a conjectured hardness assumption does not meet it.
\par\smallskip\textit{Formulation and concept references:} \sref{30006683}{1}.
\cataloguescope{Prove unconditionally that bounded-error quantum polynomial time (BQP) strictly contains bounded-error probabilistic polynomial time (BPP).}
\subsection{Short English statement}
Can one prove, without computational assumptions, that quantum computers efficiently solve some decision problem beyond all efficient randomized classical computers?
%% BEGIN RESEARCH INSERTION
\subsection{Research attempt}
\paragraph{Current frontier.}
\textit{Research check: 27 September 2026.}
Shor supplies polynomial-time quantum factoring, not a classical lower bound.
\eref{30006683}{1}. The 2025 sampling/function paper in the original bibliography
obtains conditional consequences involving collapse of the polynomial
hierarchy. The September 2026 Fourier-hierarchy manuscript concerns oracle
separations. Neither statement supplies the unconditional total language
required here. \sref{30006683}{2}, \sref{30006683}{3}. We distinguish these carefully
before trying to convert a quantum algorithm into a separating language.

\paragraph{Attack route.}
There are three distinct missing steps: prove classical hardness of a known
quantum task; make an oracle lower bound survive an explicit description of
the oracle; or turn a promise separation into a total-language separation.
The last two steps are often concealed by a change of encoding. We pursue
the first through an exact decision/search reduction, then test the other
two against that reduction. All claims of hardness below remain hypotheses.

\paragraph{Attempt: a total factoring threshold language.}
Fix a conventional, polynomial-time parsable binary encoding of pairs and
put
\[
 L_{\rm fac}=\{\langle N,t\rangle:N\geq2,\ 2\leq t<N,
                  \ \exists d\ (2\leq d\leq t,\ d\mid N)\}.
 \tag{19.1}
\]
Malformed encodings, primes, and out-of-range thresholds are no-instances.
The language is total. A bounded-error quantum factorization algorithm,
followed by multiplication checks on its output, decides it in quantum
polynomial time. Complete factorization reveals the least prime factor;
on a prime it reveals no proper factor. The usual amplification and
polynomial-time truncation of the factoring algorithm give the required
worst-case bounded-error time bound. This uses the established quantum
algorithm, not an assumption about its classical complexity. \eref{30006683}{1}.

Here is the converse reduction with adaptive error accounted for. Suppose
$L_{\rm fac}\in\mathsf{BPP}$, and let $N$ be composite of binary length $n$.
The predicate in (19.1), as a function of $t\in[2,N-1]$, changes from zero
to one precisely at the least prime divisor $p$ of $N$. Binary search uses
at most $n$ decision calls and returns $p$ if all answers are correct.
Amplify every call to error at most $1/(12(n+1))$, using fresh random bits.
Even though later queries depend on earlier answers, conditional on any
past transcript the next query is a fixed input, so its error bound still
holds. The union bound makes the entire search succeed with probability
at least $11/12$. Division verifies the output; a failed verification can
be reported as failure. The running time is polynomial in $n$, including
amplification and arithmetic. Consequently
\[
 \begin{split}
 &\text{no randomized polynomial-time algorithm finds a proper factor}\\[-2pt]
 &\text{of every composite with success at least }2/3
 \quad\Longrightarrow\quad
 L_{\rm fac}\in\mathsf{BQP}\setminus\mathsf{BPP}.
 \end{split}
 \tag{19.2}
\]
This is a complete conditional implication, not an unconditional
separation. In particular, a negative answer to UnsolvedMath problem 8400004's
\emph{deterministic} factoring question would not supply the stronger
randomized lower bound assumed here.

\paragraph{Attempt: a quantitative oracle-to-input audit.}
Consider an oracle $f:\{0,1\}^m\to\{0,1\}^m$ in Simon's promise problem:
$f(x)=f(y)$ exactly when $y=x$ or $y=x\oplus s$, for a nonzero hidden $s$.
The quantum algorithm obtains random linear equations orthogonal to $s$;
classical queries require a collision to expose it in the black-box model.
The established query-complexity framework, including the distinction
between oracle and ordinary machines, is essential here. \eref{30006683}{3}.

If the \emph{entire table} of $f$ is the input, its length is
$B=m2^m$. A classical algorithm can read and sort the pairs $(f(x),x)$ in
time polynomial in $B$, recover a collision, and obtain $s$. Thus an
$\exp(\Omega(m))$ query lower bound is not a superpolynomial lower bound
in this encoding. A proof that charged time in $m$ while giving the table
as an input would have changed the target's size parameter.

If instead a size-$\operatorname{poly}(m)$ circuit $C$ describes $f$, quantum
queries can be implemented by reversible evaluation of $C$, with its
workspace uncomputed. But a classical algorithm now sees $C$, not just
an oracle. As a concrete stress test, take a rank-$m-1$ binary matrix $A$
and $f(x)=Ax$. Its kernel is $\{0,s\}$, so it satisfies Simon's promise.
Given the coefficient list defining $C$, Gaussian elimination finds $s$
in polynomial time. The oracle difficulty therefore does not hold
uniformly for small descriptions, even though their truth tables have the
same two-to-one pattern. A hard explicit family, together with a lower
bound against algorithms inspecting its descriptions, is still required.

\paragraph{Attempt: what a legitimate totalization needs.}
Suppose a promise problem has a polynomial-time decidable promise set $D$
and a bounded-error quantum algorithm on $D$. Define a total language by
rejecting inputs outside $D$. The preliminary membership test for $D$
then gives a BQP algorithm for the totalization, and any BPP decider for
it solves the original promise problem. This elementary lemma is sound.
It is useless without a cheap test for the \emph{entire} promise.

For an arbitrary circuit, ``acceptance probability lies outside
$(1/3,2/3)$'' is a semantic promise, not a syntactic property that may be
assumed efficiently decidable. Repeating the circuit does not decide
membership in that promise with a fixed gap near either endpoint.
Likewise, testing Simon periodicity by evaluating every pair uses an
exponential number of evaluations in a succinct encoding. Rejecting
invalid descriptions is therefore an uncharged operation in the naive
construction. The total threshold language (19.1) avoids that particular
problem, but does not avoid the classical-hardness problem.

\paragraph{Stress test.}
The reduction handles all composite inputs and all adaptive transcripts;
it is not an average-case factoring assertion. Its hypothesis is stronger
than merely excluding deterministic factoring. Oracle results and sampling
results have not been substituted for a lower bound for (19.1). The
linear Simon example refutes only a universal description-to-oracle
transfer principle, not the existence of other hard descriptions.
Also, $\mathsf P=\mathsf{PSPACE}$ would imply
$\mathsf{BQP}=\mathsf{BPP}=\mathsf P$, by the standard polynomial-space
simulation of quantum computation; the target would then be false.
\eref{30006683}{2}. That dependency is another reason an unconditional
hardness assertion here cannot be supplied by a query bound alone.

\paragraph{Outcome.}
\textbf{Outcome: Conditional reduction.}
The notebook gives a total-language witness conditional on worst-case
randomized factoring hardness, with a fully charged adaptive reduction.
It also isolates concrete failures of table encoding, succinct-oracle
transfer, and unchecked promise totalization. No randomized factoring
lower bound has been established. The conditional implication and the
special examples are not claimed to be historically new.
%% END RESEARCH INSERTION
\subsection{Sources}
\begin{itemize}\raggedright
\item[\textnormal{[S1]}]\hypertarget{source:30006683:1}{} E. Bernstein, U. Vazirani, SIAM J. Comput. 26(5):1411-1473, 1997.
\url{https://doi.org/10.1137/S0097539796300921}.
{\small\textit{Catalogue formulation source.}}
\item[\textnormal{[S2]}]\hypertarget{source:30006683:2}{} Improved Separation Between Quantum and Classical Computers for Sampling and Functional Tasks — Marshall, Aaronson and Dunjko, CCC 2025.
\url{https://drops.dagstuhl.de/storage/00lipics/lipics-vol339-ccc2025/LIPIcs.CCC.2025.5/LIPIcs.CCC.2025.5.pdf}.
{\small\textit{Catalogue research/status reference; not a proof endorsement.}}
\item[\textnormal{[S3]}]\hypertarget{source:30006683:3}{} Oracle Separations in the Fourier Hierarchy — Atul Mantri, September 10, 2026.
\url{https://arxiv.org/abs/2609.11830}.
{\small\textit{Catalogue research/status reference; not a proof endorsement.}}
%% BEGIN RESEARCH REFERENCES
\item[\textnormal{[E1]}]\hypertarget{extra:30006683:1}{} Peter W. Shor, \emph{Polynomial-Time Algorithms for Prime Factorization and Discrete Logarithms on a Quantum Computer}, SIAM Journal on Computing 26(5) (1997), 1484--1509; arXiv:quant-ph/9508027.
\url{https://arxiv.org/abs/quant-ph/9508027}.
{\small\textit{Added research source. Quantum factoring algorithm; no classical hardness theorem is inferred. Consulted 27 September 2026.}}
\item[\textnormal{[E2]}]\hypertarget{extra:30006683:2}{} John Watrous, \emph{Quantum Computational Complexity}, arXiv:0804.3401v1 (2008), Sections III.2 and IV.
\url{https://arxiv.org/pdf/0804.3401v1}.
{\small\textit{Added research source. Fixed-gate conventions and polynomial-space classical simulation. Consulted 27 September 2026.}}
\item[\textnormal{[E3]}]\hypertarget{extra:30006683:3}{} Daniel R. Simon, \emph{On the Power of Quantum Computation}, SIAM Journal on Computing 26(5) (1997), 1474--1483, DOI 10.1137/S0097539796298637.
\url{https://www.cs.miami.edu/home/burt/learning/csc595.251/s0097539796298637.pdf}.
{\small\textit{Added research source. Primary paper for the hidden-XOR-period oracle problem, not an ordinary language-class separation. Consulted 27 September 2026.}}
%% END RESEARCH REFERENCES
\end{itemize}

%% END ORIGINAL SECTION CONTAINER
\end{document}
